\documentclass[10pt,a4paper]{article}
\usepackage{amsmath,amssymb,amsthm,graphicx}
% Minimal retained preamble for lane-p2 cards taken from
% sources/ahl-16-fixed-points-2019/editorial-preamble.tex
% (doi:10.5802/ahl.16, Annales Henri Lebesgue 2 (2019) 149-185, CC BY 4.0).
%
% Copied verbatim from the editorial preamble: the package list, the
% theorem declarations under the shared counter and the macro block, so that
% every macro and environment of the retained spans is available.
%
% Omitted and why: the class line (the assembler supplies the class); the
% a4paper geometry line and the fontspec block with its explicit lmroman10
% files (encodings and fonts fight the forced xelatex engine, and the
% assembler supplies the page setup).
\usepackage{amsmath,amssymb,amsthm,mathtools,graphicx,xcolor,hyperref,xurl,xspace,ifthen}
\newtheorem{theo}{Theorem}[section]
\newtheorem{lemm}[theo]{Lemma}
\newtheorem{theorem}[theo]{Theorem}
\newtheorem*{remas*}{Remarks}
\providecommand{\og}{«\nobreakspace}
\providecommand{\fg}{\nobreakspace»}
\newtheorem{lem}[theo]{Lemma}
\newtheorem{lemma}[theo]{Lemma}
\newtheorem{prop}[theo]{Proposition}
\newtheorem{coro}[theo]{Corollary}
\newtheorem{corol}[theo]{Corollary}
\newtheorem{conj}[theo]{Conjecture}
\newtheorem{Proposition}[theo]{Proposition}
\newtheorem{theoreme}[theo]{Theorem}
\theoremstyle{definition}
\newtheorem{defi}[theo]{Definition}
\newtheorem{exam}[theo]{Example}
\newtheorem{exem}[theo]{Example}
\newtheorem{rema}[theo]{Remark}
\newtheorem{nota}[theo]{Notation}
\newtheorem*{rema*}{Remark}
\newtheorem*{defi*}{Definition}
\newtheorem*{theo*}{Theorem}
\newtheorem*{lemm*}{Lemma}
\newtheorem*{prop*}{Proposition}
\newcommand{\enoncename}{}
\newtheorem{corpusenonce}[theo]{\enoncename}
\newenvironment{enonce}[2][]{\renewcommand{\enoncename}{#2}\begin{corpusenonce}}{\end{corpusenonce}}
\newenvironment{enonce*}[2][]{\par\medskip\noindent\textbf{#2.}\itshape}{\par\medskip}
\newenvironment{proofc}[1][\proofname]{\begin{proof}[#1]}{\end{proof}}
\newcommand{\Mc}{\mathcal{M}}
\newcommand{\pointir}{.\ ---\ }
\providecommand{\mathof}[1]{\mathrm{#1}}

\numberwithin{equation}{section}
\numberwithin{figure}{section}
\numberwithin{table}{section}
\def\endappendix{}
%~Mouliné par MaN_auto v.0.21.4 2019-04-30 09:02:09


\usepackage[all]{xy}

\hyphenation{auto-morphism
every
Roelcke
peri-od-ic
peri-od-icy
rep-re-sen-ta-tion
}








\newcommand{\A}{\mathcal A}
\newcommand{\N}{\mathbb{N}}
\newcommand{\Z}{\mathbb{Z}}
\newcommand{\Q}{\mathbb{Q}}
\newcommand{\K}{\mathcal{K}}
\newcommand{\CC}{\mathbb{C}}

\newcommand*{\mA}{\mathbf{A}}
\newcommand*{\mB}{\mathbf{B}}
\newcommand*{\mC}{\mathbf{C}}
\newcommand*{\mD}{\mathbf{D}}
\newcommand*{\mE}{\mathbf{E}}
\newcommand*{\mF}{\mathbf{F}}
\newcommand*{\mH}{\mathbf{H}}
\newcommand*{\mZ}{\mathbf{Z}}
%\newcommand*{\mA}{\mathbf{A}}
%\newcommand*{\mA}{\mathbf{A}}

\newcommand{\restrict}[2]{#1 \upharpoonright #2}

\newcommand{\e}{\varepsilon}

\DeclareMathOperator{\Aut}{Aut}
\DeclareMathOperator{\Stab}{Stab}
\DeclareMathOperator{\Age}{Age}

\newcommand{\actson}{\curvearrowright}

\newcommand{\RUCB}{\mathrm{RUC}_{b}}
\newcommand{\ROB}{\mathrm{Ro}_{b}}
\newcommand{\WAP}{\mathrm{WAP}}
\newcommand{\AP}{\mathrm{AP}}

\graphicspath{{./figures/}}

\newcommand*{\mk}{\mkern -1mu}
\newcommand*{\Mk}{\mkern -2mu}
\newcommand*{\mK}{\mkern 1mu}
\newcommand*{\MK}{\mkern 2mu}

\hypersetup{urlcolor=purple, linkcolor=blue, citecolor=red}

\newcommand*{\romanenumi}{\renewcommand*{\theenumi}{\roman{enumi}}}
\newcommand*{\Romanenumi}{\renewcommand*{\theenumi}{\Roman{enumi}}}
\newcommand*{\alphenumi}{\renewcommand*{\theenumi}{\alph{enumi}}}
\newcommand*{\Alphenumi}{\renewcommand*{\theenumi}{\Alph{enumi}}}
\let\oldexists\exists
\renewcommand*{\exists}{\mathrel{\oldexists}}
\let\oldforall\forall
\renewcommand*{\forall}{\mathrel{\oldforall}}


\usepackage{cleveref}
\crefname{originalsection}{original section}{original sections}
\Crefname{originalsection}{Original section}{Original sections}
\makeatletter
\newlabel{section:oscillationstability}{{3}{source}{}{sourcepointer-1}{}}
\newlabel{section:oscillationstability@cref}{{[originalsection][0][]3}{[1][0][]source}{}{}{}}
\makeatother
\title{The ambit of the identity of the Gelfand algebra generated by a uniformly continuous bounded function is isomorphic to the orbit closure of the function}
\author{Lionel Nguyen Van Thé}
\date{doi:10.5802/ahl.16; CC BY 4.0}
\begin{document}
\maketitle
\noindent\textit{Editorial note.} The counted unit is the article's own Proposition prop:Gelfand at source lines 421-423 together with its proof at source lines 434-448, a straight adjacent pairing in the source file. The source is the publisher version of the article in Annales Henri Lebesgue 2 (2019), pages 149-185, DOI 10.5802/ahl.16, released under CC BY 4.0.\par\medskip
\noindent\textit{Editorial note.} The licence evidence used here is the observation file kept with the source, license-observation.json, which records the licence as CC BY 4.0, points at the inherited publisher-version evidence and states that the landing page displays no licence line; the underlying LICENSE\_EVIDENCE.json and inherited-license.txt are the same records the earlier cards of this source rely on. No acceptance record is written anywhere.\par\medskip
\noindent\textit{Editorial note.} The retained proof is headed with the author's own title, given as the optional argument of the proof environment and naming this very proposition by reference; that label is defined inside the retained statement, so no pointer is needed for it, while the reference to the article's earlier section on oscillation stability is declared as a pointer because that section is not retained.\par\medskip
\noindent\textit{Editorial note.} The retained preamble is the source's editorial preamble with three omissions only: the class line, which the assembler supplies; the page geometry; and the fontspec block with its explicit font files, which would fight the forced xelatex engine. Everything else, including the theorem declarations under the shared counter and the whole macro block, is retained verbatim; the two omitted decorations are the only text of the abstract not reproduced.\par\medskip
\noindent\textit{Editorial note.} Neither the retained statement nor the retained proof contains a citation, so this card needs no bibliography entry.\par\medskip
\noindent\textit{Editorial note.} The census package listed the statement span as 421-423 and the proof span as 434-448 and marked the candidate free of missing citation keys; both spans and the pointer list were re-read from the source before the card was written.\par\medskip
\noindent\textit{Editorial note.} rqual sv-mismatch batch fix 2026-10-14 (Chief ruling: re-assemble unpublished only): source\_version aligned verbatim to the source dir license-observation.json sv (package evidence authority); revision bumped for re-assembly. Recipe from lane-s3.\par\medskip
\noindent\textit{Allowed standard prerequisites (editorial).}\begin{itemize}
\item Uniformly continuous bounded functions, Gelfand algebras, ambits and orbit closures. For a topological group the right-uniformly continuous bounded functions form a C-star algebra on which the group acts by left translation, and for a function in that algebra the closed unital left-invariant subalgebra it generates has a Gelfand spectrum, a compact space on which the group acts with the identity character as a distinguished point; the pair consisting of this compact space and the point is the ambit of the identity. Dually, the closure of the orbit of the function under the action is a compact space with the function itself as distinguished point. The author's own definitions and the surrounding material of the article are the ones used by the counted statement and proof.. The counted proposition identifies the two descriptions of the universal ambit attached to a function, and the article uses it to transfer statements about the Gelfand compactification to statements about orbit closures, which is the route its fixed-point results take.
\end{itemize}\medskip
\ifdefined\hypertarget\hypertarget{sourcepointer-1}{}\fi
\noindent\textit{Source reference.} Original section 3: sources/ahl-16-fixed-points-2019/body-01.tex, source lines 185--185. Navigation pointer invoked by the retained proof when it recalls the material of the article's earlier section on oscillation stability, where the function is identified with an evaluation function on the spectrum of the algebra it generates. That section is the third section of the article, so the number recorded is 3; the section is not retained in this card, the cited material is an allowed input and none of it is reproduced inside the span. The first revision recorded the label string as the number and the assembler refused it, since a reference number is required to match a numeric pattern.\par\medskip
\begin{prop}\label{prop:Gelfand}
Let $G$ be a topological group, $f\in \RUCB(G)$. Let $\langle f \rangle$ denote the unital left-invariant, closed C$^{*}$-subalgebra of $\RUCB(G)$ generated by $f$. Then the ambits $(G^{\langle f \rangle},e_{G})$ and $(\overline{G\bullet f},f)$ are isomorphic.
\end{prop}

\begin{proof}[Proof of Proposition~\ref{prop:Gelfand}]
As we have seen in Section~\ref{section:oscillationstability}, $f$ can be thought of as the continuous function $\hat f$ on $G^{\langle f \rangle}$ defined by $\hat f:\phi\mapsto \phi(f)$. It follows that for every $\phi\in G^{\langle f \rangle}$, the map $\pi(\phi):h\mapsto \hat f (h\cdot \phi) \ (=h\cdot \phi(f) = \phi(h^{-1}\cdot f))$ is in $\RUCB(G)$. This defines
\[
\pi:G^{\langle f \rangle}\longrightarrow \RUCB(G).
\]

Note that for $g,h\in G$, $\pi(\hat g)(h)=\hat g(h^{-1}\cdot f)=(h^{-1}\cdot f)(g)=f(hg)=g\bullet f(h)$. Therefore, $\pi(\hat g)=g\bullet f$, and in particular $\pi(e_{G})=f$. Let us now verify that $\pi$ is an injective homomorphism of $G$-flows. This will suffice to prove the desired result, since $\pi$ will then be a $G$-flow isomorphism between $G^{\langle f \rangle}$ and its image in $ \RUCB(G)$, which is $\overline{G\bullet \pi(e_{G})}=\overline{G\bullet f}$.

For injectivity, assume that $\pi(\phi_{1})=\pi(\phi_{2})$. From the expression of $\pi(\phi)(h)$ above, this implies that $\phi_{1}$ and $\phi_{2}$ agree on the orbit $G\cdot f$, and therefore on all of $\langle f \rangle$. To prove that $\pi$ is $G$-equivariant, consider $g,h\in G$ and $\phi\in G^{\langle f \rangle}$. Then:
\[
\pi(g\cdot \phi)(h)=(g\cdot \phi)(h^{-1}\cdot f)=\phi(g^{-1}\cdot(h^{-1}\cdot f))=\phi((hg)^{-1}\cdot f)=\pi(\phi)(hg).
\]

The last term of the equality is $(g\bullet\pi(\phi))(h)$, so $\pi(g\cdot \phi)=g\bullet\pi(\phi)$. To prove that $\pi$ is continuous, fix $H\subset G$ finite, $\e>0$. If $\phi_{1}, \phi_{2}\in G^{\langle f \rangle}$ agree on the finite set $H^{-1}\cdot f$, then $|\phi_{1}(h^{-1}\cdot f)-\phi_{2}(h^{-1}\cdot f)|<\e$ for every $h\in H$. This means that for every $h\in H$, $|\pi(\phi_{1})(h)-\pi(\phi_{2})(h)|<\e$, as required.
\end{proof}

\end{document}
